\subsection{应用：Lebesgue 积分中的变量替换}

设 I 是 R 的一个区间（有界或否），a 是它在 \(\overline{R}\) 中的左端点，b 是右端点，\(\mu\) 是 I 上的 Lebesgue 测度。对每个 \(\mu\) -可积函数 f 以及每个区间 \(H \subset I\) （左端点为 \(\alpha\) ，右端点为 \(\beta\) ），我们把 \(\int_{\alpha}^{\beta} \mathbf{f}(t) dt\) 写出来以代替 \(\int_{H} \mathbf{f}(t) dt = \int_{H} \mathbf{f} d\mu\) ，并置 \(\int_{\beta}^{\alpha} \mathbf{f}(t) dt = -\int_{\alpha}^{\beta} \mathbf{f}(t) dt\) ；当 f 是具有紧支集的正则函数时（Ch. IV, §4, No.~4, 例），这样赋予这些记号的意义与 FRV, II, §§1 和 2 中赋予它们的意义相一致。

设 \(g\) 是定义在 \(I\) 上且局部 \(\mu\) -可积的数值函数，\(x_0\) 是 \(I\) 的一个点；对每个 \(x \in I\) ，置

\[
G (x) = c + \int_ {x _ {0}} ^ {x} g (t) d t \quad (c \text {   a   constant }).\tag{11}
\]

数值函数 G 在 I 上连续；这由 Lebesgue 定理（Ch. IV, §4, No.~3, 定理 2 的推论 1）立即得出，因为 g 与以 x 和 \(x + h\) 为端点的区间的特征函数之积当 h 趋于 0 时趋于一个可忽略函数。因此 G(I) 是 R 的一个区间。在本 No. 中，我们把 G 看作一个由 I 到局部紧空间 G(I) 上的映射。用 \(\lambda\) 表示 I 上的测度 \(g \cdot \mu\) 。

先设 g 是 \(\mu\) -可积的。于是，与上面同样的推理表明极限 \(\mathrm{G}(a+)\) 和 \(\mathrm{G}(b-)\) 存在且有限；此外，测度 \(|\lambda|\) 是有界的（\(\S5\) , No.~3, 定理 1 的推论），且由 I 到 \(\mathrm{G}(\mathrm{I})\) 中的映射 G 是 \(\lambda\) -逆紧的。

命题 8. --- 设 g 是 \(\mu\) -可积的。若 J 表示 R 中以 \(\mathrm{G}(a+)\) 和 \(\mathrm{G}(b-)\) 为端点的开区间，则测度 \(g \cdot \mu\) 在 G 下的像是测度 \(\varphi_{J} \cdot \nu\) （若 \(\mathrm{G}(a+) \leqslant \mathrm{G}(b-)\) ）或是测度 \(-\varphi_{J} \cdot \nu\) （若 \(\mathrm{G}(a+) \geqslant \mathrm{G}(b-)\) ）（其中 \(\nu\) 表示 \(\mathrm{G}(\mathrm{I})\) 上的 Lebesgue 测度）。

只须证明，对每个函数 \(f \in \mathcal{K}\big(\mathrm{G}(\mathrm{I})\big)\) ，

\[
\int_ {\mathrm{G} (a +)} ^ {\mathrm{G} (b -)} f (\xi) d \xi = \int_ {a} ^ {b} f (\mathrm{G} (t)) g (t) d t.\tag{12}
\]

现在，这个公式对 \(g \in \mathcal{K}(\mathrm{I})\) 已经证明过了（FRV, II, §2, No.~1, 公式 (1)）。我们过渡到一般情形；存在一个序列 \((g_n)\) （由 \(\mathcal{K}(\mathrm{I})\) 中的函数组成），使得：\(1^\circ\) 序列 \((g_n(t))\) 在 I 中几乎处处趋于 \(g(t)\) ；\(2^\circ\) 存在一个 \(\mu\) -可积函数 \(h \geqslant 0\) ，使得 \(|g_n| \leqslant h\) 对一切 \(n\) 成立（Ch. IV, §3, No.~4, 定理 3）。由 Lebesgue 定理立即推出：置 \(G_n(x) = c + \int_{x_0}^x g_n(t) dt\) ，序列 \((G_n)\) 在 I 上一致收敛于 G，且数 \(G_n(a+)\) 和 \(G_n(b-)\) 分别趋于 \(G(a+)\) 和 \(G(b-)\) 。设 \(f'\) 是 \(\mathcal{K}(\mathbf{R})\) 中的一个延拓 \(f\) 的函数；上面的事实表明 \(f'(G_n(t))\) 趋于 \(f'(G(t)) = f(G(t))\) ，对一切 \(t \in \mathrm{I}\) 而言；应用 Lebesgue 定理即见公式 (12) 由公式

\[
\int_ {\mathrm{G} _ {n} (a +)} ^ {\mathrm{G} _ {n} (b -)} f ^ {\prime} (\xi) d \xi = \int_ {a} ^ {b} f ^ {\prime} \bigl (\mathrm{G} _ {n} (t) \bigr) g _ {n} (t) d t
\]

经取极限而得。

推论. --- 若一个定义在 G(I) 上、取值于 \(\overline{R}\) 或一个 Banach 空间中的函数 f 使得函数 \(t \mapsto \mathbf{f}(G(t))g(t)\) 在 I 上对 Lebesgue 测度可积，则 f 在 J 上对 Lebesgue 测度可积，且

\[
\int_ {\mathrm{G} (a +)} ^ {\mathrm{G} (b -)} \mathbf {f} (\xi) d \xi = \int_ {a} ^ {b} \mathbf {f} (\mathrm{G} (t)) g (t) d t\tag{13}
\]

（Lebesgue 积分中的变量替换公式）。

因为，\(\mathbf{f}\big(\mathrm{G}(t)\big)\) 对测度 \(|g|\cdot\mu\) 可积，从而对测度 \(g^{+}\cdot\mu\) 和 \(g^{-}\cdot\mu\) 也可积；由（No.~2 的）定理 1 推出，f 对像测度 \(\mathrm{G}(g^{+}\cdot\mu)\) 和 \(\mathrm{G}(g^{-}\cdot\mu)\) 可积，从而对测度 \(\varphi_{J}\cdot\nu\) 也可积，且 (13) 成立，这要注意到命题 8 和公式 (9)。

可能发生这样的情形：\(\mathbf{f}\) 在 J 上对 Lebesgue 测度可积，而 \(t \mapsto \mathbf{f}(\mathbf{G}(t))g(t)\) 在 I 上对 Lebesgue 测度不可积（习题 10）。

现在设 g 几乎处处保持常号（且是局部 \(\mu\) -可积的）；例如可以设 \(g(t) \geqslant 0\) 在 I 中几乎处处成立。于是 G 是 I 上的一个递增连续函数，因此 \(\mathrm{G}(a+)\) 和 \(\mathrm{G}(b-)\) 存在（但可以是无穷）。此外，G 是一个 \(\lambda\) -逆紧映射（由 I 到 \(\mathrm{G}(\mathrm{I})\) 中的）：因为，若 \(\mathrm{G}(b-) \in \mathrm{G}(\mathrm{I})\) ，则存在一个 \(x_{1} \geqslant x_{0}\) ，使得 G 对 \(x \geqslant x_{1}\) 是常值的，于是紧区间 \([G(x_{0}), G(b-)]\) 在 G 下的逆像是 \(\lambda\) -可积的；反之，若 \(G(b-) \notin G(I)\) ，则每个以 \(G(x_{0})\) 为左端点、含于 G(I) 中的紧区间在 G 下的逆像与一个紧区间至多相差一个 \(\lambda\) -可忽略区间。对以 \(G(x_{0})\) 为右端点的紧区间可类似地论证，由此得到我们的断言。此外：

命题 9. --- 设 \(g \geqslant 0\) 且局部 \(\mu\) -可积。则 \(G\) 作用于正测度 \(g \cdot \mu\) 所得的像是 \(G(I)\) 上的 Lebesgue 测度。为使一个函数 \(f\) （定义在 \(G(I)\) 上、取值于 \(\overline{R}\) 或一个 Banach 空间中）在 \(G(I)\) 上对 Lebesgue 测度可积，必须且只须函数 \(t \mapsto f(G(t))g(t)\) 在 \(I\) 上对 Lebesgue 测度可积，此时关系式 (13) 成立。

陈述的第一部分由下列事实得出：公式 (12) 对每个函数 \(f \in \mathcal{K}(\mathrm{G}(\mathrm{I}))\) 成立；因为，依上面的注记，函数 \(t \mapsto f(\mathrm{G}(t))\) 的支集含于一个区间 \(K \subset I\) ，g 在其上是可积的，只须对 K 应用命题 8。第二部分是 No.~2 的定理 1 的推论。

